\section{问题重述与总体思路}
在半径 $1800\,\mathrm m$ 的圆域内，有若干位置未知、频道互异的干扰源。机器狗从圆心出发，逐次选择检测位置和频道，根据带有 $\pm1^\circ$ 误差的示向度推断目标位置，并在距离目标不超过 $20\,\mathrm m$ 时完成光学定位与清除。任务要求同时控制移动、切换频道、检测及清除产生的时间\cite{problem,protocol}。

问题一要求给出交会多边形的直径算法，并判断相应直径圆能否覆盖区域。问题二要求在已有一次示向读数后选择第二检测点。问题三要求在全向源、数量未知的条件下尽快完成全部清除。问题四加入定向源，本文保留解法位置。

前三问具有递进关系：问题一提供位置不确定性的几何表达；问题二解决信息获取的位置选择；问题三将单目标定位嵌入多目标覆盖和路径规划。模型采用集合估计描述必然性，用附加概率模型比较典型效率，用执行日志评价实际耗时。
\placeholder{总体模型流程}{绘制四个模块：覆盖扫描、按频道更新可行域、选择目标并补测、光学清除；补测模块返回可行域更新，清除后返回目标调度；最右侧为逐频道终止证书。第四问不绘制。}

\section{基本假设与适用范围}
\begin{enumerate}
\item 干扰源位置、频道和有效接收半径在一次任务中保持不变；不同源的频道互异且信号互不影响。这是将多目标观测按频道分解的依据。
\item 采用题设二维平面、直线移动与恒速模型。不另加障碍物、转弯、加减速及仪器姿态调整代价；检测耗时已包含题设天线调整过程。
\item 示向误差严格处于 $[-1^\circ,1^\circ]$，同一地点同一频道的误差固定。确定性分析不要求误差服从某种分布，也不要求不同地点误差独立。
\item 观测与清除反馈遵守题面和通信协议。全向源在接收半径内返回示向或近距反馈；光学清除只由距离决定。收到成功反馈后才将该频道标记为已清除。
\item 为比较问题二候选点的平均性能，额外采用源位置在目标圆内等面积均匀、接收半径在 $[1000,1500]$ 上均匀、不同检测位置角误差独立均匀的参考先验。这些是评价假设，不能作为题面事实或清除保证。
\item 数值计算采用向外包络和容差复核。空集合、退化多边形及浮点计算异常触发检查或重新检测，不能直接解释为目标不存在。
\end{enumerate}
目标数量在10至16之间只能用于一致性检查。清除10个目标并不能证明任务结束；全向源的全域覆盖证书才可排除未发现目标。

\section{符号说明}
\begin{longtable}{>{$}l<{$}p{9.5cm}}
\toprule
\multicolumn{1}{l}{符号}&含义及单位\\\midrule\endhead
\Omega=\disk(0,R_0)&目标圆域，$R_0=1800\,\mathrm m$；$\disk$ 表示闭圆盘\\
G_j,c_j,\rho_j&第 $j$ 个源的位置、频道及固定有效接收半径，$1000\le\rho_j\le1500\,\mathrm m$\\
a,b&接收半径下、上界，分别为1000和1500米\\
S_i,\theta_i&第 $i$ 个检测点与示向度；推导中角度统一用弧度\\
\eps&最大角误差，$\pi/180$\\
e(\theta),e_\perp(\theta)&$(\cos\theta,\sin\theta)$ 及其逆时针旋转 $90^\circ$ 的单位向量\\
W_i,K_j,\widehat K_j&示向楔形、源位置可行集及其保守多边形外包络\\
D(K),C(K),r_*(K)&区域直径、最小包围圆圆心和半径\\
r_{\rm opt},v&光学作用半径 $20\,\mathrm m$，移动速度 $5\,\mathrm{m/s}$\\
\mathcal P_j,\mathcal N_j&频道 $j$ 的有效接收记录与无信号记录集合\\
\mathcal V,\mathcal C&第二检测点的稳健接收区域与进一步筛选的候选区域\\
r_h,h&覆盖扫描环半径；兜底网格边长上界\\
L,N_m,N_s,N_o,N_c&总路程、检测次数、实际切频次数、光学尝试次数、成功清除数\\
T,\bar t,\eta&虚拟总时间、每源平均定位清除时间、清除比例\\
\bottomrule
\end{longtable}
为简化下标，讨论单个源时省略频道编号。二维叉积定义为 $u\times w=u_xw_y-u_yw_x$，$\operatorname{wrap}$ 将角度映射至 $(-\pi,\pi]$。
